% LTeX: enabled=false
% ==============================================================================
% SHORTCUT COMMANDS
% ==============================================================================
% Provides \DefineScName and \DefineScNames to quickly define commands that
% render their content in \textsc. The command name is derived automatically
% from the text (e.g. "Lorem Ipsum" becomes \LoremIpsum) or can be specified
% explicitly.
% ==============================================================================

\ExplSyntaxOn

% -----------------------------------------------------------------------
% Internal message
% -----------------------------------------------------------------------

% tssuite sc-command-already-defined: raised when trying to define a
%   small-caps command whose name is already taken.
\msg_new:nnn{tssuite}{sc-command-already-defined}
  {SC~command~'#1'~is~already~defined.}

% -----------------------------------------------------------------------
% Internal variables
% -----------------------------------------------------------------------

% Accumulates the derived command name from concatenated words.
\tl_new:N \l__tssuite_sc_cmd_tl

% Holds the full original text to display in \textsc.
\tl_new:N \l__tssuite_sc_text_tl

% Holds the space-split parts of the input text.
\seq_new:N \l__tssuite_sc_parts_seq

% -----------------------------------------------------------------------
% Internal functions
% -----------------------------------------------------------------------

% \tssuite_define_sc_command:Nn <cs> <text>
%   Define <cs> as a parameterless command that prints <text> in \textsc.
%   #1 - control sequence token (e.g. \loremipsum)
%   #2 - text to render in small caps
\cs_new_protected:Npn \tssuite_define_sc_command:Nn #1 #2 {
  \cs_if_exist:NTF #1 {
    \msg_error:nnn {tssuite} {sc-command-already-defined} {\token_to_str:N #1}
  } {
    \NewDocumentCommand #1 { } {\textsc{#2}}
  }
}

% \tssuite_define_sc_command:cn <cs name> <text>
%   Like \tssuite_define_sc_command:Nn but takes the command name as a string.
%   #1 - command name as a string (e.g. "loremipsum")
%   #2 - text to render in small caps
\cs_new_protected:Npn \tssuite_define_sc_command:cn #1 #2 {
  \cs_if_exist:cTF {#1} {
    \msg_error:nnn {tssuite} {sc-command-already-defined} {#1}
  } {
    \exp_args:Nc \NewDocumentCommand {#1} { } {\textsc{#2}}
  }
}

% \tssuite_make_sc_name_command:n <text>
%   Derive a command name from <text> by concatenating non-blank words,
%   then define that command to render <text> in \textsc.
%   #1 - full text (e.g. "Lorem Ipsum" → \LoremIpsum)
\cs_new_protected:Npn \tssuite_make_sc_name_command:n #1 {
  \tl_set:Nn \l__tssuite_sc_text_tl {#1}
  \tl_trim_spaces:N \l__tssuite_sc_text_tl
  \seq_set_split:NnV \l__tssuite_sc_parts_seq {~} \l__tssuite_sc_text_tl
  \tl_clear:N \l__tssuite_sc_cmd_tl
  \seq_map_inline:Nn \l__tssuite_sc_parts_seq {
    \tl_if_blank:nF {##1} {
      \tl_put_right:Nn \l__tssuite_sc_cmd_tl {##1}
    }
  }
  \exp_args:NVV \tssuite_define_sc_command:cn
    \l__tssuite_sc_cmd_tl
    \l__tssuite_sc_text_tl
}

% -----------------------------------------------------------------------
% User commands
% -----------------------------------------------------------------------

% \DefineScName[<cmd>]{<text>}
%   Define a new small-caps command.
%   If the optional <cmd> is omitted, the command name is derived from <text>
%   by concatenating its space-separated words (e.g. "Foo Bar" → \FooBar).
%   #1 - (optional) explicit command to define
%   #2 - text to render in \textsc
\NewDocumentCommand \DefineScName { o m } {
  \IfNoValueTF {#1}
    {\tssuite_make_sc_name_command:n {#2}}
    {\tssuite_define_sc_command:Nn #1 {#2}}
}

% \DefineScNames{<csv>}
%   Define multiple small-caps commands from a comma-separated list of texts.
%   Each entry follows the auto-naming rule of \DefineScName (no optional arg).
%   #1 - comma-separated list of texts, e.g. {Foo Bar, Baz Qux}
\NewDocumentCommand \DefineScNames { >{\SplitList{,}} m } {
  \ProcessList {#1} {\tssuite_make_sc_name_command:n}
}

\ExplSyntaxOff
